\section{Introduzione all'Aggregate Computing}
\label{sec:introduzione-aggregate-computing}
%
L'aggregate computing è un paradigma di programmazione che affronta la complessità dello sviluppo di sistemi distribuiti e decentralizzati astraendo la computazione su una rete di dispositivi interagenti come se fosse un'unica entità collettiva. L'elemento concettuale alla base del paradigma è il campo computazionale, una funzione che associa a ogni punto dello spazio e del tempo occupato dai dispositivi un valore, manipolabile attraverso composizione funzionale. Operativamente, ogni dispositivo esegue periodicamente, in modo asincrono rispetto agli altri, un round computazionale: percepisce lo stato dei propri vicini e dell'ambiente, esegue localmente lo stesso programma aggregato e ne comunica l'esito, contribuendo così, attraverso l'iterazione di questo ciclo nel tempo, all'emergere del comportamento collettivo desiderato a partire da sole interazioni locali. Questa astrazione consente di specificare in modo dichiarativo il comportamento desiderato dell'intero sistema, con meccanismi di coordinamento intrinsecamente robusti e scalabili al variare del numero di dispositivi partecipanti.

Dal punto di vista formale, queste astrazioni trovano fondamento nel field calculus, un calcolo che fornisce i costrutti di base -- diffusione dell'informazione, raccolta e aggregazione dei dati, evoluzione nel tempo -- a partire dai quali è possibile costruire funzioni di livello più alto. Su queste basi sono state sviluppate librerie di comportamenti collettivi auto-stabilizzanti, sia di natura generica sia specializzati per domini applicativi specifici, tra cui la robotica di sciame e il coordinamento gerarchico: è proprio in questo filone che si inserisce \textit{MacroSwarm} \cite{aguzzi2024macroswarm}, la libreria su cui si basa il progetto descritto in questa tesi.

Negli anni sono stati sviluppati diversi framework a supporto dell'aggregate computing, ciascuno dei quali implementa i costrutti fondamentali del field calculus offrendo al contempo un proprio approccio all'esecuzione e all'integrazione con l'ambiente circostante. Tra questi, \textit{ScaFi} \cite{casadei2022scafi} si distingue come ambiente di sviluppo basato su Scala, pensato per essere utilizzabile sia in simulazione sia in scenari di esecuzione reale, e offre un'API di alto livello per la definizione concisa di computazioni aggregate, mantenendo al contempo l'estensibilità necessaria per l'integrazione con diversi ambienti di esecuzione. Un altro framework rilevante è \textit{FCPP} \cite{audrito2024fcpp}, le cui implementazioni leggere ne consentono l'esecuzione anche su microcontrollori e schede embedded a basso consumo, rendendo il paradigma applicabile anche a reti IoT e a sistemi robotici con risorse computazionali limitate, sebbene nel progetto descritto in questa tesi la computazione aggregata sia demandata a un nodo esterno anziché eseguita direttamente a bordo dei singoli robot, come illustrato nella sezione seguente.